\section{Introduction}

\bd[Derivative]
A \textbf{derivative} is a contract for a future transacion between two parties, that derives its value from the
performance of an underlying entity.
\ed

\bd[Underlying]
The \textbf{underlying} is an asset, index, interest rate, or any variable from which a derivative is derived from.
\ed

Derivatives can be either ``exchange-traded'' or ``over-the-counter''.

\bd[Exchange-Traded Derivative]
An \textbf{exchange-traded derivative} is a standardized derivative that is traded on a regulated derivative exchange.
\ed

\bd[Derivative Exchange]
A \textbf{derivative exchange} is a central financial exchange where people can trade standardized exchange-traded
derivatives defined by the exchange.
\ed

Once two traders have agreed to a trade offered by an exchange, it is handled by the clearing house.

\bd[Clearing House]
A \textbf{clearing house} is an organization that facilitates tradings between two parties and manages the risks.
\ed

The advantage of a clearing house is that traders do not have to worry about the creditworthiness of the people they
are trading with. The clearing house takes care of credit risk by requiring each of the two traders to deposit funds
to ensure that they will live up to their obligations. \v

Traditionally exchange-traded derivatives traders used to physically meet on the floor of the exchange, shouting, and
using a complicated set of hand signals to indicate the trades they would like to carry out. Exchanges have largely
replaced this system by electronic trading which involves the use of algorithms to initiate trades, often without human
intervention, and has become an important feature of derivatives markets.

\bd[Over-The-Counter Derivative]
An \textbf{over-the-counter (OTC) derivative} is a derivative that is traded off-exchange.
\ed

Banks, other large financial institutions, fund managers, and corporations are the main participants in OTC derivatives
markets. \v

Once an OTC trade has been agreed, the two parties can either present it to a central counterparty (CCP), which is like
an exchange clearing house, or clear the trade bilaterally where the two parties have usually signed an agreement
covering all their transactions with each other.

\bd[Central Counterparty]
A \textbf{central counterparty (CCP)} stands between the two parties to the derivatives transaction so that one party
does not have to bear the risk that the other party will default.
\ed

While OTC markets were largely unregulated, following the financial crisis of 2008 we have seen the development of
many new regulations affecting the operation of OTC markets. The main objectives of the regulations are to improve
the transparency of OTC markets and reduce systemic risk. The over-the-counter market in some respects is being forced
to become more like the exchange-traded market. \v

Both OTC and exchange-traded markets for derivatives are huge. The number of derivatives transactions per year in OTC
markets is smaller than in exchange-traded markets, but the average size of the transactions is much greater, making the
volume of business in OTC market much larger than in the exchange-traded one. \v

Derivatives markets have been outstandingly successful. The main reason is that they have attracted many different 
types of traders and have a great deal of liquidity. When a trader wants to take one side of a contract, 
there is usually no problem in finding someone who is prepared to take the other side. \v

Three broad categories of traders can be identified: ``hedgers'', ``speculators'', and ``arbitrageurs''. 

\bd[Hedger]
A \textbf{hedger} is a trader who uses derivatives to reduce the risk that they face from potential future movements in
a market variable.
\ed

\bd[Speculator]
A \textbf{speculator} is a trader who uses derivatives to bet on the future direction of a market variable.
\ed

\bd[Arbitrageur]
An \textbf{arbitrageur} is a trader who takes offsetting positions in two or more instruments to lock in a profit.
\ed

The fact that derivatives can be used for hedging, speculation, and arbitrage, makes them very versatile instruments.
Sometimes traders who have a mandate to hedge risks or follow an arbitrage strategy become (consciously or
unconsciously) speculators. The results can be disastrous, so it is very important for both financial and
nonfinancial corporations to set up controls to ensure that derivatives are being used for their intended purpose.
Risk limits should be set and the activities of traders should be monitored daily to ensure that these risk limits
are adhered to.

Unfortunately, even when traders follow the risk limits that have been specified, big mistakes can happen. Some of
the activities of traders in the derivatives market during the period leading up to the start of the financial crisis
in July 2007 proved to be much riskier than they were thought to be by the financial institutions they worked for.

\section{Forwards \& Futures}

\bd[Forward / Future]
A \textbf{forward} or \textbf{future} is a contract between two parties to buy or sell an underlying asset at a
specified price at a specified future date.
\ed

Forwards are traded in OTC markets, usually between two financial institutions or between a financial institution and
one of its clients. Futures are traded on exchanges, and the exchange acts as the counterparty to all trades.

\bd[Long Position]
The party that agrees to buy the asset is said to be in a \textbf{long position}.
\ed

\bd[Short Position]
The party that agrees to sell the asset is said to be in a \textbf{short position}.
\ed

\bd[Delivery Date / Maturity Date]
The specified future date $T$ in a forward or future is called the \textbf{delivery date} or \textbf{maturity date}.
\ed

\bd[Delivery Price / Forward Price / Future Price]
The specified price $K$ in a forward or future is called the \textbf{delivery price}, or \textbf{forward price}, or
\textbf{future price}.
\ed

Indepentendly of the delivery price, the underlying asset has its own price at each time step, called the ``spot
price''.

\bd[Spot Price]
The \textbf{spot price} $S_t$ is the price of an asset for immediate delivery at time t.
\ed

The payoff from a long position in a forward or future on one unit of an asset is:
\bse
S_T - K
\ese

where $K$ is the delivery price and $S_T$ is the spot price of the asset at maturity of the contract. This is because
the holder of the contract is obligated to buy an asset worth $S_T$ for $K$.

Similarly, the payoff from a short position in a forward of future contract on one unit of an asset is:
\bse
K - S_T
\ese

These payoffs can be positive or negative. Because it costs nothing to enter into a forward or future contract, the
payoff from the contract is also the trader's total gain or loss from the contract.

\vspace{5pt}
\fig{dv1}{0.65}

\section{Options}

\bd[Option]
An \textbf{option} is a contract that gives the holder the right, but not the obligation, to buy or sell an underlying
asset at a specified price at a specified future date.
\ed

It should be emphasized that an option gives the holder the right to do something. The holder does not have to
exercise this right. This is what distinguishes options from forwards and futures, where the holder is obligated to
buy or sell the underlying asset.

\bd[Strike Price / Exercise Price]
The specified price $K$ in an option is called the \textbf{strike price} or \textbf{exercise price}.
\ed

\bd[Expiration Date / Maturity]
The specified future date $T$ in an option is called the \textbf{expiration date} or \textbf{maturity}.
\ed

Options are traded both on exchanges and in the over-the-counter market. Based on if the option can be exercised only
on the expiration date or at any time up to the expiration date, we distinguish between ``European'' and ``American''
options.

\bd[European Option]
A \textbf{European option} is an option that can be exercised only on the expiration date.
\ed

\bd[American Option]
An \textbf{American option} is an option that can be exercised at any time up to the expiration date.
\ed

Note that the terms American and European do not refer to the location of the option or the exchange. Some options
trading on North American exchanges are European. Most of the options that are traded on exchanges are American.
European options are generally easier to analyze than American options, and some of the properties of an American
option are frequently deduced from those of its European counterpart. \v

There are two types of options: ``calls'' and ``puts''.

\bd[Call Option]
A \textbf{call option} is an option that gives the holder the right to buy an underlying asset at a specified price
at a specified future date.
\ed

\bd[Put Option]
A \textbf{put option} is an option that gives the holder the right to sell an underlying asset at a specified price
at a specified future date.
\ed

At this stage we note that there are four types of participants in options markets:
\bit
\item Buyers of calls.
\item Sellers of calls.
\item Buyers of puts.
\item Sellers of puts.
\eit

As with forwards and futures, buyers are referred to as having long positions and sellers are referred to as having
short positions. Selling an option is also known as ``writing'' the option. \v

Whereas it costs nothing to enter into a forward or futures contract, there is a cost to acquiring an option.]

\section{Time Value Of Money}

\bd[Time Value Of Money]
The \textbf{time value of money} is the idea that money available at the present time is worth more than the same
amount in the future due to its potential earning capacity.
\ed

The reason why time value of money is valid is because of the existence of interests.

\bd[Interest]
\textbf{Interest} is the cost of borrowing money, or equivalently, the income from lending money.
\ed

\bd[Interest Rate]
The \textbf{interest rate} $r$ is the amount of interest due per period, as a proportion of the amount lent, deposited
or borrowed.
\ed

Due to interest rates, money available at the present time is worth more than the same amount in the future. This is
because money deposited in a savings account will earn interest and therefore will be worth more in the future. As a
result of this, money at present carry a ``future value''.

\bd[Future Value]
The \textbf{future value} $FV$ is the value of a current amount of money at a future date based on an assumed interest
rate of growth.
\ed

In simple words, future value is what an amount of money today will be worth in the future. This is because of the
interest that the money can earn over time, therefore making it more valuable in the future. \v

Using the interest rate one can compute the future value of an amount of money invested today. Before doing so, we
need to distinguish between simple and compound interest.

\bd[Simple Interest]
\textbf{Simple interest} is the interest on a loan or deposit that is calculated only on the amount originally
invested or borrowed.
\ed

The future value $FV$ of an amount $PV$ invested today at a simple interest rate $r$ per period for $n$ periods is:
\bse
FV = PV \cdot (1 + rn)
\ese

\bd[Compound Interest]
\textbf{Compound interest} is the interest calculated on the initial principal and also on the accumulated interest
of previous periods of a deposit or loan.
\ed

The future value $FV$ of an amount $PV$ invested today at a compound interest rate $r$ per period for $n$ periods is:
\bse
FV = PV \cdot (1 + r)^n
\ese

The difference is that with compound interest the interest for each period is added to the principal at the start of
that period, and the interest for the next period is calculated on the new principal. What this means is that
interest is being earned on both the initial investment, and the interest earned from the previous periods. \v

Future value can be contrasted with present value.

\bd[Present Value]
The \textbf{present value} $PV$ is the current value of a future sum of money or stream of cash flows given a specified
discount rate.
\ed

In simple words, future cash flows are discounted at the discount rate, and the higher the discount rate, the lower
the present value of the future cash flows. \v

The present value $PV$ of an amount $FV$ to be received in $n$ periods from now given a simple discount rate $r$ is:
\bse
PV = \frac{FV}{1 + rn}
\ese










\bd[Net Present Value (NPV)]
The \textbf{net present value} $NPV$ of a series of cash flows $CF_1, CF_2, \ldots, CF_n$ occurring at times
$t_1, t_2, \ldots, t_n$ is the sum of the present values of the individual cash flows:
\bse
NPV = \sum_{i=1}^n PV(CF_i)
\ese
\ed





